\documentclass[25pt, a0paper, portrait]{tikzposter}

\usepackage{fontspec}
\setsansfont{Helvetica Neue}
\setmainfont{Helvetica Neue}
\renewcommand{\familydefault}{\sfdefault}
\usepackage{unicode-math}
\setmathfont{FiraMath-Regular.otf}[Path=/Users/linjingxian/Library/Fonts/]
\usepackage{booktabs}
\usepackage{tikz}
\usetikzlibrary{shapes.geometric, arrows.meta, positioning, calc}
\usepackage{enumitem}
\usepackage{caption}
\usepackage[
  backend=biber,
  style=apa,
  citestyle=apa,
  natbib=false,
  doi=false,
  isbn=false,
  url=false,
  maxcitenames=2
]{biblatex}
\addbibresource{../shared/delivery_emission.bib}

% ---- Color Theme ----
\definecolorstyle{HKUStyle}{
  \definecolor{colorOne}{HTML}{1B5E20}
  \definecolor{colorTwo}{HTML}{E8F5E9}
  \definecolor{colorThree}{HTML}{2E7D32}
}{
  % Background
  \colorlet{backgroundcolor}{white}
  \colorlet{framecolor}{colorOne}
  % Title
  \colorlet{titlefgcolor}{white}
  \colorlet{titlebgcolor}{colorOne}
  % Block
  \colorlet{blocktitlebgcolor}{colorThree}
  \colorlet{blocktitlefgcolor}{white}
  \colorlet{blockbodybgcolor}{white}
  \colorlet{blockbodyfgcolor}{black}
  % Inner block
  \colorlet{innerblocktitlebgcolor}{colorTwo}
  \colorlet{innerblocktitlefgcolor}{colorOne}
  \colorlet{innerblockbodybgcolor}{white}
  \colorlet{innerblockbodyfgcolor}{black}
}

\usetheme{Default}
\usecolorstyle{HKUStyle}
\useblockstyle{Default}

\tikzposterlatexaffectionproofoff

\makeatletter
\settitle{
  \color{titlefgcolor}
  \begin{minipage}[c]{0.12\linewidth}
    \includegraphics[width=\linewidth]{hku_logo_only.pdf}
  \end{minipage}\hfill
  \begin{minipage}[c]{0.85\linewidth}
    \centering
    {\fontsize{72}{86}\selectfont\bfseries
      Scaling Emissions Assessment of Drone and\\[0.15em]
    E-Bike Food Delivery: A Gravity-Based Framework\par}
    \vspace*{1em}{\LARGE \@author \par}
    \vspace*{0.5em}{\Large \@institute}
  \end{minipage}
}
\makeatother

\title{Scaling Emissions Assessment of Drone and E-Bike Food Delivery: A Gravity-Based Framework}
\author{Presenter: Jingxian Lin \quad Supervisor: Dr.\ Jintao Ke}
\institute{Department of Civil Engineering, The University of Hong Kong}

\begin{document}

\maketitle[width=0.95\textwidth]

\begin{columns}

  % ==================== LEFT COLUMN ====================
  \column{0.5}

  \block{Background}{
    \textbf{Drone delivery: an emerging industry}
    \begin{itemize}[leftmargin=1.5em, itemsep=0.1em, topsep=0.1em]
      \item Drones are a \textbf{transformative technology} for transportation and logistics
      \item \textbf{Low-altitude economy} gaining momentum in China, included in \textbf{China's Government Work Report} in 2024 and 2025
      \item Global drone delivery market: \$3.47B (2025) $\rightarrow$ \$20.98B (2034), CAGR \textbf{19.45\%} \parencite{fortunebi2026drone}
    \end{itemize}

    \vspace{0.5em}
    \textbf{Why food delivery?}
    \begin{itemize}[leftmargin=1.5em, itemsep=0.1em, topsep=0.1em]
      \item \textbf{Current constraints:} battery capacity and cargo weight limit drones to short-distance, last-mile logistics
      \item Drone offers \textbf{food delivery advantages}: lighter packages, shorter distances ($<$ 3\,km), high time sensitivity
      \item China: \textbf{545 million} users, market growing from $\sim$300B to over \textbf{1 trillion Yuan} (2017--2027) \parencite{statista2024chinafooddelivery}
    \end{itemize}

    \vspace{0.5em}
    \textbf{Challenges for drone deployment}
    \begin{itemize}[leftmargin=1.5em, itemsep=0.1em, topsep=0.1em]
      \item \textbf{High costs} of purchasing and maintaining drones; significant \textbf{UAS infrastructure} investment required
      \item \textbf{Lack of predictive models} for environmental impact; unclear if drones offer a \textbf{sustainable advantage}
    \end{itemize}
    \vspace{0.5em}
    \begin{center}
      \begin{tikzpicture}
        \node[draw=colorThree, fill=colorTwo, rounded corners=10pt, inner sep=14pt, text width=0.96\linewidth, align=left, line width=3pt] {
          \fontsize{38}{58}\selectfont
          \textbf{Key Gap:} Most existing studies \textbf{lack real origin-destination (OD) data} and rely on aggregate metrics that cannot capture route-level emission variations.\par
        };
      \end{tikzpicture}
    \end{center}
  }

  \block{Research Objectives}{
    \begin{enumerate}[leftmargin=1.5em, itemsep=0.1em, topsep=0.1em]
      \item Develop a \textbf{gravity-based demand generation method} producing door-to-door order data for precise emission analysis.
      \item Conduct the \textbf{first comprehensive comparison} of carbon emissions between e-bike and drone delivery across multiple Chinese cities.
      \item Identify \textbf{key urban features} that determine when drone delivery achieves lower emissions.
    \end{enumerate}
  }

  \block{Interim Results}{
    \begin{minipage}[c]{0.48\linewidth}
      \textbf{Emission Reduction Potential}
      \begin{itemize}[leftmargin=1.5em, itemsep=0.1em]
        \item Threshold-based assignment: orders above distance $d^*$ go to drones, below to e-bikes. Even with e-bike bundling ($\kappa{=}0.7$--$0.9$), optimal drone deployment achieves significant emission savings.
      \end{itemize}

      \vspace{0.5em}
      \textbf{Key Findings}
      \begin{itemize}[leftmargin=1.5em, itemsep=0.1em]
        \item All-drone delivery reduces emissions by \textbf{33\%} on average vs.\ all e-bike; optimal drone--e-bike mix achieves \textbf{36\%} mean reduction.
        \item \textbf{Route factor} (detour ratio) is the strongest predictor of drone emission advantage, with benefits increasing alongside road network complexity and delivery distance.
        \item Results are robust across \textbf{seven} diverse urban districts in China.
      \end{itemize}
    \end{minipage}\hfill
    \begin{minipage}[c]{0.48\linewidth}
      \begin{center}
        \includegraphics[width=\linewidth]{../../figures/district/nanshan/nanshan_emission_reduction.pdf}
      \end{center}
    \end{minipage}
  }

  \block{Future Work}{
    \begin{itemize}[leftmargin=1.5em, itemsep=0.1em]
      \item Incorporate \textbf{weather impacts} (wind, rain) on drone energy consumption.
      \item Extend to other delivery modes (e.g., gasoline cars, electric cars, autonomous vehicles, etc.).
      \item Develop a \textbf{predictive model} for optimal drone deployment ratios across urban typologies.
    \end{itemize}
  }

  % ==================== RIGHT COLUMN ====================
  \column{0.5}

  \block{Methods}{
    \vspace{0.3em}
    \begin{center}
      \begin{tikzpicture}[scale=2.0,
          data/.style={rectangle, draw=blue!60, fill=blue!10, minimum width=3cm, minimum height=1.2cm, align=center, font=\normalsize},
          output/.style={rectangle, draw=gray!70, fill=gray!10, rounded corners=5pt, minimum width=3cm, minimum height=1.2cm, align=center, font=\normalsize},
          smallbox/.style={rectangle, draw=gray!90, fill=gray!90, text=white, minimum width=3cm, minimum height=0.7cm, align=center, font=\small, rounded corners=3pt},
          model/.style={diamond, draw=teal!70, fill=teal!15, minimum width=1.8cm, minimum height=1.5cm, align=center, font=\normalsize, inner sep=2pt},
          arrow/.style={-{Stealth[scale=1.0]}, very thick, draw=black!80}
        ]
        % --- Column 1: Input pair ---
        \node[data] (sales) at (0, 0) {Monthly Sales};
        \node[data] (building) at (0, -2) {Building\\Footprint};
        % --- Column 2: Gravity Model ---
        \node[model] (gravity) at (3.2, -1) {Gravity\\Model};
        % --- Column 3: OD Matrix + Road Network (same y as sales & building) ---
        \node[output] (od) at (6.4, 0) {OD Matrix};
        \node[data] (roadnet) at (6.4, -2) {Road\\Network};
        % --- Column 4: Shortest Path (same y as Gravity Model) ---
        \node[model] (shortest) at (9.6, -1) {Shortest\\Path};
        % --- Column 5: Result ---
        \node[output] (result) at (12.8, -1) {OD Matrix};
        \node[smallbox, below=0cm of result.south, anchor=north] (airdist) {Air Distance};
        \node[smallbox, below=0cm of airdist.south, anchor=north] (roaddist) {Road Distance};
        % --- Column 6: Analysis ---
        \node[rectangle, draw=colorOne, fill=colorTwo, rounded corners=5pt, minimum width=3cm, minimum height=1.2cm, align=center, font=\normalsize, right=1.5cm of result] (analysis) {Analysis};
        % --- Arrows ---
        \draw[arrow] (sales.east) -- (gravity);
        \draw[arrow] (building.east) -- (gravity);
        \draw[arrow] (gravity) -- (od);
        \draw[arrow] (od.east) -- (shortest);
        \draw[arrow] (roadnet.east) -- (shortest);
        \draw[arrow] (shortest) -- (result);
        \draw[arrow] (result) -- (analysis);
      \end{tikzpicture}
    \end{center}

    \vspace{0.5em}
    % --- Step 1: Data & OD Generation ---
    \textbf{Step 1: Generating realistic delivery orders.}\enspace
    Real food-delivery platforms do not publish individual order origins and destinations. We reconstruct them using a \textbf{gravity model}---a well-established spatial interaction model that distributes trips between locations based on their ``attractiveness'' and mutual distance.

    \vspace{0.2em}
    \textit{Input data.}\enspace Monthly sales records of \textbf{1.07 million merchants} provide store locations and order volumes. Building footprints from mapping APIs give potential customer locations.

    \vspace{0.2em}
    \textit{Gravity model.}\enspace We assume that larger buildings attract more orders due to higher population density, while greater distances reduce demand owing to increased delivery time and cost. Each merchant $i$ with demand $P_i$ distributes orders to nearby buildings $j$ in proportion to building volume $v_j$ and inversely with distance $d_{ij}$:
    \[
      D_{ij} = P_{i} \cdot \frac{v_{j}^\alpha \,/\, d_{ij}^{\beta}}{\sum_{k \in \mathcal{D}} v_{k}^\alpha \,/\, d_{ik}^{\beta}}
    \]
    where $\alpha{=}1.7$ and $\beta{=}0.5$ are calibrated so the mean delivery distance matches the observed 2.23\,km \parencite{zhang_exploration_2023}.

    \vspace{0.2em}
    \begin{center}
      \raisebox{-0.5\height}{\includegraphics[width=0.46\linewidth]{../../figures/district/nanshan/nanshan_store_coverage.png}}\hfill
      \raisebox{-0.5\height}{\includegraphics[width=0.46\linewidth]{../../figures/remote/luohu/luohu_road_distance_distribution.pdf}}
    \end{center}

    \vspace{0.2em}
    \textbf{Step 2: Computing delivery distances.}\enspace
    For each OD pair we compute the \textbf{straight-line (air) distance} for drones and the \textbf{road-network distance} for e-bikes via shortest-path routing on OpenStreetMap. Their ratio is the \textbf{route factor}.

    \vspace{0.2em}
    \begin{center}
      \raisebox{-0.5\height}{\includegraphics[width=0.36\linewidth]{../../figures/district/nanshan/nanshan_distance_comparison.png}}\hfill
      \raisebox{-0.5\height}{\includegraphics[width=0.56\linewidth]{../../figures/district/nanshan/nanshan_distance_detour_scatter.pdf}}
    \end{center}

    \vspace{0.2em}
    % --- Step 3: Emission Models ---
    \textbf{Step 3: Lifecycle emission modelling} \parencite{kirschstein_comparison_2020, figliozzi_lifecycle_2017}\textbf{.}
    \begin{itemize}[leftmargin=1.5em, itemsep=0.1em]
      \item \textbf{Use phase:} physics-based power models compute energy per trip for each vehicle---e-bike (rolling resistance, aerodynamic drag, grade, acceleration) and drone (induced lift, blade drag, fuselage drag, climb). Energy $\times$ grid emission factor gives use-phase CO$_2$e.
      \item \textbf{Production \& disposal:} embodied emissions from manufacturing and end-of-life processing of vehicle body and batteries, amortised over the vehicle's \textbf{lifetime delivery count}.
      \item \textbf{Comparison:} Lifecycle emission comparison vs.\ delivery distance and route factor (road/air distance ratio). Green = drone produces fewer emissions. Drones gain advantage at longer distances and higher route factors.
    \end{itemize}

    \vspace{0.2em}
    \begin{center}
      \includegraphics[width=0.92\linewidth]{../../figures/lifecycle/lifecycle_distance_route_factor.pdf}
    \end{center}
    \vspace{-0.5em}

    % \vspace{0.5em}
    % % --- Step 5: Multi-City Comparison ---
    % \textbf{Step 5: Cross-region comparison.}\enspace
    % We apply the full pipeline to \textbf{six urban regions} across China with diverse terrain, road networks, and population densities. For each region, we compare all-e-bike, all-drone, and optimally-mixed deployment strategies.
  }

  % \block{Power Consumption Analysis}{
  %   \begin{center}
  %     \includegraphics[width=0.95\linewidth]{../../figures/sensitivity/sensitivity_velocity_power_stacked.pdf}
  %   \end{center}
  %   \vspace{-0.5em}
  %   {\small Power consumption breakdown for drone (a) and e-bike (b) at varying velocities. Induced power dominates drone consumption; rolling resistance dominates e-bike consumption.}
  % }

  % \block{Lifecycle Sensitivity Analysis}{
  %   \begin{center}
  %     \includegraphics[width=0.92\linewidth]{../../figures/lifecycle/lifecycle_distance_route_factor.pdf}
  %   \end{center}
  %   \vspace{-0.5em}
  %   {\small Lifecycle emission comparison vs.\ delivery distance and route factor (road/air distance ratio). Green = drone produces fewer emissions. Drones gain advantage at longer distances and higher route factors.}
  % }

  % \block{Study Regions}{
  %   \begin{center}
  %     \includegraphics[width=0.8\linewidth]{../../figures/analysis/regions_map.pdf}
  %   \end{center}
  %   \vspace{-0.5em}
  %   {\small Six study regions across China spanning diverse urban forms, terrain, and climate conditions.}
  % }

  % \block{Cross-Region Comparison}{
  %   \begin{center}
  %     \includegraphics[width=0.95\linewidth]{../../figures/analysis/annual_emissions_comparison.pdf}
  %   \end{center}
  %   \vspace{-0.5em}
  %   {\small Annual lifecycle emissions across six regions. Optimal mixed drone--e-bike deployment achieves \textbf{23--48\%} emission reduction compared to all e-bike delivery.}
  % }

  % \block{Emission Reduction Potential}{
  %   \begin{center}
  %     \includegraphics[width=0.88\linewidth]{../../figures/district/nanshan/nanshan_emission_reduction.pdf}
  %   \end{center}
  %   \vspace{-0.5em}
  %   {\small Threshold-based assignment: orders above distance $d^*$ go to drones, below to e-bikes. Even with e-bike bundling ($\kappa{=}0.6$--$0.9$), optimal drone deployment achieves significant emission savings.}
  % }

  % \block{Key Findings \& Impact}{
  %   \begin{tikzpicture}
  %     \node[draw=colorThree, fill=colorTwo, rounded corners=10pt, inner sep=15pt, text width=0.92\linewidth, align=left] {
  %       \begin{itemize}[leftmargin=1.2em, itemsep=0.3em]
  %         \item All-drone delivery reduces emissions by \textbf{27\%} vs.\ all e-bike; optimal drone--e-bike mix achieves \textbf{30\%} reduction.
  %         \item \textbf{Route factor} (road detour) is the strongest predictor of drone emission advantage, with benefits increasing alongside \textbf{road network complexity} and \textbf{delivery distance}. Results are robust across diverse urban environments.
  %         \item \textbf{Impact:} Provides actionable distance- and route-factor-based criteria for city planners and logistics operators to determine where drone deployment yields the greatest emission savings.
  %       \end{itemize}
  %     };
  %   \end{tikzpicture}
  % }

  % \block{Future Work}{
  %   \begin{itemize}[leftmargin=1.5em, itemsep=0.2em]
  %     \item Incorporate \textbf{weather impacts} (wind, rain) on drone energy consumption.
  %     \item Integrate \textbf{economic viability} analysis combining emission and cost models.
  %     \item Extend to additional delivery modes: \textbf{autonomous vehicles} and \textbf{truck-drone hybrids}.
  %     \item Develop a \textbf{predictive model} for optimal drone deployment ratios across urban typologies.
  %   \end{itemize}
  % }

  % \block{References}{
  %   {\footnotesize
  %     \printbibliography[heading=none]
  %   }
  % }

\end{columns}

\end{document}
